\begin{afterword}
\summaryhead

The post draws three lessons from Asch's conformity experiments: a unanimous group can make many people say black is white, a single dissenter sharply reduces conformity, and conformity returns when the dissenter sides with the group. Dissent helps a group but costs the dissenter, who must keep it up and may be wrong.

An anecdote shows how groups become overconfident. When the author raised possible problems with a project, a newcomer answered each with reassurance. Alone, people doubt themselves; together, politeness stops them from questioning each other, and each takes the other's reassurance for confidence. At the scale of a group this is ``pluralistic ignorance.''

The author and Robin Hanson agree that common knowledge of a factual disagreement means someone is irrational, so for rationalists disagreement is serious. The most important lesson, the post says, is to distinguish expressing a concern from disagreeing. But socially there is little difference: speaking out cannot be undone, and the group will not thank you. The decision is left to the reader.

\responsehead

The post names a real problem well. Doubts that each person keeps to themselves can add up to a group that is surer than any of its members, and a single doubt said aloud can lead others to voice theirs. The problems are its premise about disagreement and its ending. Its summary of Asch, that dissent ``was not learned,'' is Asch's own first reading of one condition; the notes on ``Asch's Conformity Experiment'' give the others.

The theoretical premise is stated without its contested step. Aumann's theorem assumes that the parties share a prior, that is, start from the same beliefs before any evidence, and Hanson's own statements of the result say so. To conclude that common knowledge of a disagreement shows ``someone must be irrational,'' one needs the further claim that rational agents must share priors. Hanson has argued for a version of that claim, and Hanson's paper notes that others reject the usual argument for it. The post gives neither the premise nor the dispute.

The post is candid about the limits of its lesson. It calls the distinction between a concern and a disagreement ``the most important lesson,'' and two paragraphs later says that socially there is not much difference, that people are not wired for it, and that a group of rationalists could only ``pretend.'' The last paragraph leaves the decision to the reader, with no guidance on when a concern is worth its cost.

In short: a sound observation about suppressed doubts, a premise about disagreement stated without its contested step, and a lesson the post itself says does not hold socially, left without guidance on when to speak.
\end{afterword}
